\begin{center}
\begin{minipage}{\linewidth}

\captionof{table}{Ablation results on the independent diagnostic target
pool. Gain is the paired MSE reduction relative to the reference candidate.}
\label{tab:ablation}

\centering
\scriptsize
\renewcommand{\arraystretch}{1.08}
\setlength{\tabcolsep}{3.0pt}

\begin{tabular*}{\linewidth}{@{\extracolsep{\fill}}lccccc@{}}
\toprule

Variant
& MSE$\downarrow$
& Worst-10\%$\downarrow$
& Gain (\%)$\uparrow$
& \shortstack{Harmful selection\\(\%)$\downarrow$}
& \shortstack{Limit satisfaction\\(\%)$\uparrow$}
\\

\midrule

\textbf{RCF-DTI}
& \textbf{0.004063}
& \textbf{0.015764}
& \textbf{15.79}
& \textbf{15.00}
& \textbf{100.00}
\\

\shortstack[l]{Earlier pooled-source training\\for non-reference models}
& 0.004876
& 0.019044
& 1.75
& 2.50
& 100.00
\\

Reference candidate only
& 0.005051
& 0.019512
& 0.00
& 0.00
& 100.00
\\

Two reference-architecture initializations
& 0.004917
& 0.019136
& 1.29
& 1.25
& 100.00
\\

Without reference regularization$^{\dagger}$
& 0.004055
& 0.015690
& 16.16
& 17.50
& 100.00
\\

Without complexity filtering$^{\dagger}$
& 0.004037
& 0.015635
& 16.13
& 13.75
& 96.25
\\

\bottomrule
\end{tabular*}

\vspace{2pt}
\footnotesize

$^{\dagger}$ Diagnostic variants only.
The harmful-selection rate is computed over all diagnostic cases.
Without complexity filtering, a selected model may exceed an assigned
limit.

\end{minipage}
\end{center}
